% =====================================================================
% Introdução: Rivera-Alvarez & Bejan (2003)
%
% Citações no padrão abnTeX2 (pacote abntex2cite):
%   \cite{chave}        ->  (AUTOR, ano)
%   \citeonline{chave}  ->  Autor (ano)
% Se o seu documento usa natbib, troque \cite por \citep
% e \citeonline por \citet.
% =====================================================================
\section{Introdução}
\label{sec:introducao}

A separação de misturas gasosas é uma etapa central em muitos processos
industriais e ambientais. Alguns exemplos são a produção de nitrogênio e
oxigênio a partir do ar, a purificação de hidrogênio, a desumidificação de
correntes gasosas, a recuperação de hidrocarbonetos e, mais recentemente, a
captura de dióxido de carbono. Entre as tecnologias disponíveis, a adsorção se
destaca por usar sólidos porosos que retêm preferencialmente um componente (ou
uma família de componentes) da alimentação. Essa seletividade pode vir do
equilíbrio, quando as afinidades termodinâmicas são diferentes, ou da cinética,
quando as velocidades de difusão no sólido são diferentes
\cite{ruthven1984,ruthven1994}.

Como a capacidade do adsorvente é finita, todo processo de adsorção é cíclico.
Depois da etapa de adsorção vem uma etapa de regeneração, ou dessorção, em que
a espécie retida é liberada e o sólido recupera sua capacidade. A regeneração
pode ser feita pela redução da pressão (\textit{Pressure Swing Adsorption},
PSA) ou pelo aumento da temperatura (\textit{Temperature Swing Adsorption},
TSA). Esse caráter intermitente transforma o tempo em variável de projeto. Um
ciclo curto demais aproveita pouco a capacidade do sólido em cada etapa. Um
ciclo longo demais mantém o adsorvente perto da saturação, onde a taxa de
transferência de massa é baixa. Existe, portanto, um ritmo de operação ótimo,
especialmente importante nos processos controlados pela cinética.

Ao mesmo tempo, cresce a demanda por equipamentos compactos, capazes de
transferir o máximo de massa (ou de calor) por unidade de volume e de tempo.
Reduzir as escalas de comprimento acelera a difusão, mas obriga a distribuir o
fluido por um número muito grande de canais pequenos, o que aumenta a potência
de bombeamento. Projetar um trocador compacto é, portanto, buscar um
compromisso entre transferência e escoamento. A resposta a esse compromisso é
uma arquitetura de escoamento, e não apenas um conjunto de dimensões.

A teoria constructal, formulada por Adrian Bejan na segunda metade da década de
1990, oferece um método para tratar esse tipo de problema. Seu princípio
central, a lei constructal, afirma que, para um sistema de escoamento de
tamanho finito persistir no tempo, sua configuração deve evoluir de modo a
facilitar o acesso das correntes que o atravessam
\cite{bejan2000,bejan2008}. Na engenharia, essa lei se traduz num procedimento
de construção hierárquica. Primeiro, identifica-se o menor elemento capaz de
cumprir a função desejada e otimiza-se sua forma sob restrições globais. Em
seguida, esses elementos são agrupados em construtos cada vez maiores, também
otimizados \cite{bejan1997}. O resultado típico são estruturas dendríticas, em
forma de árvore, semelhantes às de bacias hidrográficas, redes vasculares e vias
aéreas pulmonares \cite{bejan2012}.

Na natureza, ritmo e estrutura aparecem juntos. Na respiração, inspiração e
expiração se alternam em intervalos de mesma ordem de grandeza. O ar é
distribuído por uma árvore de vias aéreas até as regiões onde a troca de massa
ocorre por difusão. O pulmão é, nesse sentido, um trocador de massa
intermitente com estrutura hierárquica, que é justamente o que um equipamento
de adsorção compacto procura ser.

Nesse contexto, \citeonline{riveraalvarez2003} aplicam o método constructal a
processos de adsorção--dessorção e tratam de forma integrada a organização
temporal e espacial do equipamento. Na primeira parte do trabalho, com modelos
deliberadamente simples, os autores mostram que existe um tempo de adsorção que
maximiza a taxa média de captura. Eles analisam três situações: tempo de
dessorção fixo, tempo de dessorção variável e concentração variável da espécie
adsorvida na fase gasosa. Na segunda parte, as escalas de tempo otimizadas
determinam as dimensões do volume elementar, isto é, do canal de gás e da camada
adsorvente. A partir dele são montados hierarquicamente o primeiro, o segundo e
o terceiro construtos. Em cada nível, o número de elementos e a distribuição do
volume ocupado pelos canais são escolhidos para minimizar a potência de
bombeamento. A arquitetura resultante é formada por duas árvores de escoamento
acopladas ``copa com copa'', uma de alimentação e outra de coleta. A principal
conclusão é que há oportunidades reais de otimizar o ritmo e a forma no projeto
de trocadores de massa com máxima densidade de transferência por unidade de
volume.

% ---------------------------------------------------------------------
% Parágrafo opcional (relatório próprio): remova os "%" para usar.
% ---------------------------------------------------------------------
% Neste trabalho, a metodologia proposta por \citeonline{riveraalvarez2003} é
% reproduzida analiticamente. As equações apresentadas são verificadas, os
% passos intermediários omitidos no artigo são explicitados e o procedimento é
% então aplicado a um caso com propriedades físicas reais, para avaliar as
% dimensões ótimas e a potência de bombeamento resultantes.
